%% kapa: exam=1 round=21 year=2010 period=2 form=A subject=accounting subject_ko=회계학 seq=24 number=24
%% source: Q-Net 감정평가사 기출문제 (cst00309) artlSeq=5024181 file=extracted/1008/5024181/2010년 제21회 감정평가사 1차시험 기출문제/2교시.hwp
%% conversion: kordoc 4.24.51 -> tools/gen1.py (auto); review=proofread-sonnet; flags=none
\begin{question}{24}
(주)감평은 20×1년 초 액면금액 ￦100,000의 (주)한국의 사채를 ￦95,026에 취득하고, 취득 후 상각후원가로 측정한다. (주)한국 사채의 만기는 20×3년 말이고 취득당시 유효이자율은 연 10\%, 표시이자율은 연 8\%이며 이자지급일은 매년 말이다. 20×1년 말 이자는 정상적으로 수취되었다. 그러나 20×2년 초 (주)한국의 재무상태 악화로 액면금액은 만기에 회수할 수 있지만 추후 이자수취는 불가능한 것으로 판단하여 손상차손을 인식한다면, 그 손상차손은 얼마인가? (단, 가장 근사치를 선택할 것)
\begin{center}
\small
\begin{tabular}{|c|c|c|}
\hline
\multirow{2}{*}{기간} & \multicolumn{2}{c|}{10\%} \\
\cline{2-3}
 & \makecell{1 회금액 ￦1의\\현재가치} & \makecell{기말연금 ￦1의\\현재가치} \\
\hline
1 & 0.90909 & 0.90909 \\
\hline
2 & 0.82645 & 1.73554 \\
\hline
3 & 0.75131 & 2.48685 \\
\hline
\end{tabular}
\end{center}
% [원문 그대로] 표의 '1 회금액'
\begin{choices}
  \item ￦12,381
  \item ￦12,431
  \item ￦13,884
  \item ￦13,914
  \item ￦14,015
\end{choices}
\end{question}
\answer{3}
